\documentclass[a4paper, serif, 10pt]{moderncv}

%% moderncv
% color and style
\moderncvcolor{black}
\moderncvstyle{banking}

% renew moderncv command to adapt for CJK colon
\renewcommand*{\cvitem}[3][.25em]{%
  \ifstrempty{#2}{}{\hintstyle{#2}：}{#3}%
  \par\addvspace{#1}}

\renewcommand*{\cvitemwithcomment}[4][.25em]{%
  \savebox{\cvitemwithcommentmainbox}{\ifstrempty{#2}{}{\hintstyle{#2}：}#3}%
  \setlength{\cvitemwithcommentmainlength}{\widthof{\usebox{\cvitemwithcommentmainbox}}}%
  \setlength{\cvitemwithcommentcommentlength}{\maincolumnwidth-\separatorcolumnwidth-\cvitemwithcommentmainlength}%
  \begin{minipage}[t]{\cvitemwithcommentmainlength}\usebox{\cvitemwithcommentmainbox}\end{minipage}%
  \hfill% fill of \separatorcolumnwidth
  \begin{minipage}[t]{\cvitemwithcommentcommentlength}\raggedleft\small\itshape#4\end{minipage}%
  \par\addvspace{#1}}

%% page layout/margins
\usepackage[top=2.5cm, bottom=2.5cm, left=1.5cm, right=1.5cm]{geometry}

\nopagenumbers{}

%% Babel config for English language
\usepackage[english]{babel}

%% fontspec
\usepackage{fontspec}

\IfFontExistsTF{Linux Libertine}{
  \setmainfont[Ligatures={TeX, Common}, Numbers=Lining, ItalicFont=Linux Libertine]{Linux Libertine}
}{}
\IfFontExistsTF{Linux Libertine O}{
  \setmainfont[Ligatures={TeX, Common}, Numbers=Lining, ItalicFont=Linux Libertine O]{Linux Libertine O}
}{}

%% CTeX
% CJK support, used to show CJK characters in the resume
%
% - heading=false: disable ctex heading
% - punct=kaiming: use kaiming punctuations styles for CJK
% - scheme=plain: use plain scheme, do not override \normalsize font size
% - space=auto: space settings for CJK characters
%
% CJK fonts are resolved per locale language, e.g. Noto Serif/Sans CJK JP
% for Japanese resumes, with fallback to the next region if not installed
%
% ref:
% - http://ctan.mirrorcatalogs.com/language/chinese/ctex/ctex.pdf
\usepackage[UTF8, heading=false, punct=kaiming, scheme=plain, space=auto]{ctex}

\IfFontExistsTF{Noto Serif CJK SC}{
  \setCJKmainfont{Noto Serif CJK SC}
}{}
\IfFontExistsTF{Noto Sans CJK SC}{
  \setCJKsansfont{Noto Sans CJK SC}
}{}

%% line spacing
\usepackage{setspace}
\setstretch{1.125}

%% URL styling - use normal text instead of monospace
\urlstyle{same}

%% Patch \httplink and \httpslink to handle full URLs
% The original moderncv macros always prepend http:// or https:// to URLs.
% This causes issues when the URL already has a protocol (e.g., https://example.com)
% resulting in malformed URLs like https://https://example.com.
% This patch checks if the URL already contains "://" and uses it directly if so.
% Note: We use \str_set:Nx (x-expansion) to expand macros like \@homepage first.
\ExplSyntaxOn
\str_new:N \g_yamlresume_url_str

\RenewDocumentCommand{\httpslink}{O{}m}{
  \str_set:Nx \g_yamlresume_url_str {#2}
  \str_if_in:NnTF \g_yamlresume_url_str {://}
    { \url{#2} }
    { \url{https://#2} }
}

\RenewDocumentCommand{\httplink}{O{}m}{
  \str_set:Nx \g_yamlresume_url_str {#2}
  \str_if_in:NnTF \g_yamlresume_url_str {://}
    { \url{#2} }
    { \url{http://#2} }
}
\ExplSyntaxOff

\name{陈昊然}{}
\title{人工智能工程师 | 大模型与机器学习}
\phone[mobile]{+86 138 0013 8000}
\email{chen.haoran.ai@example.com}
\homepage{https://www.haoran-ai.com}

\address{中国，北京市，北京，北京市海淀区中关村大街 27 号，100080}{}{}

\extrainfo{{\small \faGithub}\ \href{https://github.com/haoran-ai}{@haoran-ai} {} {} {} • {} {} {} 
{\small \faLinkedin}\ \href{https://www.linkedin.com/in/haoran-chen-ai}{@haoran-chen-ai} {} {} {} • {} {} {} 
{\small \faZhihu}\ \href{https://www.zhihu.com/people/haoran-ai}{@haoran-ai}}

\begin{document}

\maketitle

\section{简介}

\cvline{}{\begin{itemize}
\item 人工智能工程师，专注大语言模型、RAG 与机器学习工程化，具备从数据、训练、评估到部署的全链路经验。
\item 熟练使用 PyTorch、TensorFlow 与 Transformers，熟悉向量数据库、LangChain 和 MLOps，能够将算法原型转化为稳定服务。
\item 有推荐系统、智能问答、文本理解等落地经验，注重模型效果、推理成本与线上稳定性。
\item 具备跨团队协作与技术分享能力，持续关注生成式 AI、多模态与智能体方向。
\end{itemize}}
\section{教育背景}

\cventry{2016年9月–2019年6月}
        {硕士，计算机科学与技术，成绩：3.9/4.0}
        {清华大学}
        {\url{https://www.tsinghua.edu.cn}}
        {}
        {\begin{itemize}
\item 研究方向为自然语言处理与机器学习，参与实验室文本理解与知识抽取项目。
\item 完成大规模文本分类与实体识别实验，发表 1 篇中文会议论文。
\item 担任机器学习与数据结构课程助教，负责习题课与项目辅导。
\end{itemize}
\textbf{课程}：机器学习、深度学习、自然语言处理、计算机视觉、概率图模型、分布式系统、算法设计与分析}
\section{工作经历}

\cventry{2022年7月至今}
        {高级人工智能工程师}
        {云智科技有限公司}
        {\url{https://www.yunzhi-ai.com}}
        {}
        {\begin{itemize}
\item 负责企业级大模型应用平台的核心算法与工程架构，支持知识库问答、文档摘要和智能客服等场景。
\item 设计 RAG 检索增强生成链路，优化文档切分、召回、重排与提示词策略，问答准确率提升 23\%。
\item 主导 7B/13B 开源模型的领域微调与 LoRA 实验，建立离线评测与人工反馈闭环，降低幻觉率。
\item 推动模型服务容器化与灰度发布，结合缓存、批处理与量化将推理成本降低约 35\%。
\item 与产品、数据和后端团队协作，制定技术方案与迭代计划，指导 3 名初级算法工程师。
\end{itemize}
\textbf{关键字}：大语言模型、RAG、LoRA、PyTorch、MLOps、模型量化}

\cventry{2019年7月–2022年6月}
        {人工智能工程师}
        {星图数据科技有限公司}
        {\url{https://www.xingtu-data.com}}
        {}
        {\begin{itemize}
\item 参与个性化推荐系统研发，负责特征工程、召回与排序模型迭代，提升点击率与转化率。
\item 使用 TensorFlow 构建深度学习排序模型，设计离线训练与在线推理流程，支撑千万级用户请求。
\item 搭建 A/B 实验与模型监控看板，分析关键指标并快速定位效果波动。
\item 与数据平台团队合作优化 Spark 特征流水线，缩短模型训练数据准备时间。
\end{itemize}
\textbf{关键字}：推荐系统、TensorFlow、特征工程、Spark、A/B 实验}
\section{语言}

\cvline{中文}{母语或双语水平 \hfill \textbf{关键字}：普通话}
\cvline{英语}{完全专业水平 \hfill \textbf{关键字}：CET-6、技术文档阅读、英文技术分享}
\section{专业技能}

\cvline{机器学习与深度学习}{专家 \hfill \textbf{关键字}：PyTorch、TensorFlow、Transformers、模型微调、强化学习}
\cvline{大模型应用}{高级 \hfill \textbf{关键字}：LLM、RAG、LangChain、向量数据库、Prompt Engineering}
\cvline{工程与部署}{高级 \hfill \textbf{关键字}：Python、Docker、Kubernetes、FastAPI、MLOps}
\cvline{数据处理}{高级 \hfill \textbf{关键字}：SQL、Spark、Pandas、数据清洗}
\section{荣誉}

\cventry{2023年12月}
        {云智科技有限公司}
        {优秀员工奖}
        {}
        {}
        {\begin{itemize}
\item 因在大模型应用平台项目中的突出贡献获奖。
\item 主导 RAG 链路优化与推理成本治理，推动平台在多个业务线落地。
\end{itemize}}

\cventry{2017年10月}
        {清华大学}
        {国家奖学金}
        {}
        {}
        {\begin{itemize}
\item 研究生阶段学业成绩与科研表现优异。
\item 参与自然语言处理方向研究，完成多项实验与论文撰写工作。
\end{itemize}}
\section{证书}

\cventry{2021年8月}
        {Amazon Web Services}
        {AWS Certified Machine Learning - Specialty}
        {\url{https://aws.amazon.com/certification/certified-machine-learning-specialty/}}
        {}
        {}

\cventry{2020年3月}
        {Google}
        {TensorFlow Developer Certificate}
        {\url{https://www.tensorflow.org/certificate}}
        {}
        {}
\section{出版刊物}

\cventry{2021年10月}
        {基于预训练语言模型的中文文本分类方法研究}
        {中国计算机学会自然语言处理专委会会议}
        {\url{https://doi.org/10.12345/nlp.2021.001}}
        {}
        {\begin{itemize}
\item 研究预训练语言模型在中文文本分类任务中的迁移能力与微调策略。
\item 通过对比实验分析不同预训练模型、数据增强与提示学习方法的效果。
\item 在多个中文文本分类数据集上验证方法的有效性与稳定性。
\end{itemize}}
\section{项目}

\cventry{2023年1月–2023年12月}
        {基于 RAG 的企业知识库问答平台，支持文档解析、检索与答案生成。}
        {企业知识库智能问答系统}
        {\url{https://github.com/haoran-ai/enterprise-rag}}
        {}
        {\begin{itemize}
\item 设计文档解析、向量化、召回、重排与生成一体化链路。
\item 引入多路召回与答案引用，提升复杂问题的可追溯性与准确率。
\item 建立离线评测集与线上反馈机制，持续优化提示词与模型策略。
\end{itemize}
\textbf{关键字}：RAG、大模型、向量检索、FastAPI}

\cventry{2021年3月–2022年12月}
        {面向客服对话的意图识别、情绪分析与质检辅助平台。}
        {智能客服对话分析平台}
        {\url{https://github.com/haoran-ai/dialog-insight}}
        {}
        {\begin{itemize}
\item 构建对话文本清洗、意图分类与情绪识别流水线。
\item 使用预训练模型提升小样本场景下的分类效果。
\item 为质检团队提供可视化分析看板与高风险对话提醒。
\end{itemize}
\textbf{关键字}：NLP、意图识别、情绪分析、可视化}
\section{兴趣爱好}

\cvline{技术社区}{开源贡献、技术写作、学术分享}
\cvline{运动}{长跑、羽毛球、徒步}
\section{志愿服务}

\cventry{2020年1月–2022年12月}
        {人工智能方向志愿者}
        {中国计算机学会青年计算机科技论坛}
        {\url{https://www.ccf.org.cn}}
        {}
        {\begin{itemize}
\item 参与人工智能主题论坛与技术沙龙的筹备与组织工作。
\item 协助邀请讲者、整理活动资料并维护社区交流群。
\item 推动高校学生与产业界工程师之间的技术交流。
\end{itemize}}
    
\end{document}